\chapter{ドーパミンの本当の言葉}
% full version: chapters/07_dofaalt.tex

前の章のドーパミンは単純だった。「予想よりよいか悪いか」という1つの数である。これは第一近似であり、よく機能する。だが近年の実験は、この「配線」を1つの数より多くのものが流れていることを示した。この章は、話がもっと複雑になるところ、そして研究者がまだ議論しているところについてである。

\section*{楽観派と悲観派}

ドーパミンニューロンはどれも同じではない。うれしい驚きにいやな驚きより強く反応するものがいる。これが楽観派だ。逆のものは悲観派である。各ニューロンが自分の規則で学ぶなら、楽観派は最良の場合の報酬を、悲観派は最悪の場合の報酬を覚える。両者が一緒になると、記録されるのは1つの平均的な期待ではなく、ありうる結果のばらつき全体になる。本命だけでなく、各馬の勝ち目まで知っているブックメーカーのように。この考えは実験とよく合っている。なぜ脳がばらつき全体を必要とするのかは、まだ仮説である。

\pencilfig{7}{\pencilart{7}}{各ドーパミンニューロンは、ありうる報酬の曲線上の自分の点を受けもつ。全員でその曲線全体を覚えている。}

\section*{「よい」だけでなく「重要」も}

ドーパミンニューロンの一部は、不快だが重要な出来事、たとえば大きな音や危険にも連打で応える。ドーパミンのネットワークを流れているのは、驚きの符号だけでなくその重要度、つまり「注意せよ」という信号でもあるらしい。

\section*{1本の配線に2つの信号}

速い連打は学習させる。一方、数分かけて変わるドーパミンの遅い平均レベルは、別のことを語っているようだ。今の環境がどれだけ豊かか、である。周りに報酬が多ければ時間は貴重であり、動物はすばやく行動する。実験は、ドーパミンのレベルが行動の活発さと結びついていることを裏づけている。エンジニアならこれを周波数による信号の分離と呼ぶだろう。高い周波数はあることを、低い周波数は別のことを運ぶ。

\section*{ゴール前の上昇}

動物が報酬に近づくと、ドーパミンは連打するのではなく、なめらかに上がっていくことが多い。一つの説明は、これが期待そのもの、「ゴールは近い、がんばれ」という信号だというものだ。もう一つの説明は驚きの考えを守る。遠くからは動物は状況をぼんやりと見積もり、近づくにつれて正確になり、その一つ一つの精密化が小さなうれしい驚きになる。仮想の通路で動物を瞬時にゴールの近くへ移す巧みな実験は、後者を支持している。ドーパミンは、驚きにふさわしく連打で応えたのだ。

\section*{過去を振り返る}

大胆な代案もある。ドーパミンは「何が起こるか」ではなく「報酬の原因は何だったか」に答えている、というものだ。未来ではなく過去を見ているのである。2つの理論が異なる予測をする実験は、今のところ互いに食い違っている。これは進行中の科学論争だ。

\section*{数ではなくベクトル}

最後に、ドーパミンは価値が変わらなくても報酬の\emph{味}が変わると反応し、報酬がまったくなくても学習を助ける。どうやらこれは1つの誤差ではなく誤差の束、つまり起きたことの特徴ごとに1つずつの誤差らしい。

この章の結論。新しい理論のほとんどは単純な描像を否定するのではなく、拡張する。ドーパミンは1本の配線ではなく、数本からなる小さなケーブル束である。単純な描像と本気で対立しているのは「過去を振り返る」説だけだ。

\fullversion{7}
